\chapter{Analyse et spécification des besoins}
\label{chap:besoins}

\begin{introchapitre}
Ce chapitre traduit la problématique en exigences précises. Il identifie les acteurs du
système, puis détaille les besoins fonctionnels et non fonctionnels, avant de les représenter
par un diagramme de cas d'utilisation et la description des cas principaux.
\end{introchapitre}

\section{Identification des acteurs}

\begin{table}[htbp]
\centering
\small
\begin{tabularx}{\linewidth}{L{3.8cm}L{2.4cm}Y}
\toprule
\textbf{Acteur} & \textbf{Type} & \textbf{Rôle} \\
\midrule
Administrateur (opérateur) & humain & se connecte à l'interface, suit l'état de l'infrastructure, traite les incidents, gère les règles et les notifications \\
Agent de collecte & système & collecte les métriques, détecte les problèmes et fait évoluer les alertes ; fonctionne sans intervention humaine \\
Service de notification & système externe & Slack (webhook) et serveur SMTP qui reçoivent les messages d'alerte et de résolution \\
\bottomrule
\end{tabularx}
\caption{Acteurs du système}
\label{tab:acteurs}
\end{table}

\section{Besoins fonctionnels}

\begin{longtable}{L{1.2cm}L{4.2cm}p{8.6cm}}
\toprule
\textbf{Réf.} & \textbf{Besoin} & \textbf{Description} \\
\midrule
\endfirsthead
\toprule
\textbf{Réf.} & \textbf{Besoin} & \textbf{Description} \\
\midrule
\endhead
BF-01 & Collecter les métriques & CPU, mémoire, disque et charge de l'hôte ; CPU, mémoire, réseau et redémarrages de chaque conteneur ; séries choisies de tout exportateur Prometheus ; toutes les 5 secondes. \\
BF-02 & Gérer des règles de seuil & créer, modifier, activer, désactiver et supprimer des règles (métrique, opérateur, seuil, durée, sévérité). Les règles par défaut sont fournies dans un fichier YAML. \\
BF-03 & Détecter les anomalies & signaler une valeur anormale par rapport au comportement récent d'une métrique, sans seuil écrit à la main. \\
BF-04 & Gérer le cycle de vie des alertes & n'ouvrir une alerte qu'après la durée définie, ne jamais la dupliquer, la résoudre automatiquement quand la condition disparaît. \\
BF-05 & Regrouper en incidents & rattacher à un même incident les alertes d'une même entité apparues dans une fenêtre de temps. \\
BF-06 & Expliquer les alertes & conserver les valeurs mesurées qui justifient l'alerte et une capture des processus les plus consommateurs. \\
BF-07 & Notifier & envoyer un message Slack et un e-mail au déclenchement et à la résolution, avec nouvelles tentatives en cas d'échec et heures silencieuses. \\
BF-08 & Traiter une alerte & acquitter, retirer l'acquittement, résoudre manuellement, rouvrir. \\
BF-09 & Visualiser en temps réel & état de chaque entité, courbes avec plage attendue, incidents, mise à jour sans rechargement de la page. \\
BF-10 & Consulter l'historique & métriques sur 15 min, 1 h, 6 h et 24 h ; chronologie et post-mortem d'un incident. \\
BF-11 & Opérer les conteneurs & consulter les journaux d'un conteneur et le redémarrer depuis l'interface. \\
BF-12 & Authentifier & connexion par identifiant et mot de passe, changement obligatoire du mot de passe initial, déconnexion. \\
BF-13 & Configurer & paramétrer Slack, SMTP, heures silencieuses et détection d'anomalies depuis l'interface. \\
BF-14 & Déployer en une commande & démarrer l'ensemble et indiquer quand la plateforme est réellement prête. \\
\bottomrule
\caption{Besoins fonctionnels}
\label{tab:bf}
\end{longtable}

\section{Besoins non fonctionnels}

\begin{table}[htbp]
\centering
\small
\begin{tabularx}{\linewidth}{L{3.2cm}Y}
\toprule
\textbf{Exigence} & \textbf{Critère mesurable} \\
\midrule
Simplicité & installation par une seule commande ; aucune dépendance externe (pas de base de données à administrer). \\
Réactivité & cycle de collecte de 5 s ; interface mise à jour toutes les 2 s. \\
Fiabilité & aucune notification perdue en cas de panne d'un canal ou de redémarrage ; résolution automatique des alertes. \\
Sécurité & API fermée par défaut ; aucun mot de passe par défaut ; accès Docker minimal ; conteneurs sans privilèges. \\
Maintenabilité & tests automatisés exécutés à chaque modification ; schéma de base versionné. \\
Portabilité & tout hôte disposant de Docker (Linux, Windows, macOS ; x86-64 et ARM64). \\
Sobriété & rétention limitée (24 h brutes, 7 jours agrégées) pour maîtriser la taille de la base. \\
Ergonomie & statuts lisibles sans dépendre de la couleur ; utilisable au clavier ; affichage adapté aux petits écrans. \\
\bottomrule
\end{tabularx}
\caption{Besoins non fonctionnels}
\label{tab:bnf}
\end{table}

\section{Diagramme de cas d'utilisation global}

La figure~\ref{fig:cu-global} représente les interactions entre les acteurs et le système.

\begin{figure}[htbp]
\centering
\resizebox{\linewidth}{!}{% Diagramme de cas d'utilisation global (TikZ, notation UML)
\begin{tikzpicture}[
  font=\sffamily\footnotesize,
  uc/.style={ellipse, draw=bleuSentinel, thick, fill=bleuSentinel!6, minimum width=4.5cm, minimum height=0.85cm, align=center, inner sep=1pt},
  lk/.style={draw=grisTexte},
  inc/.style={-{Stealth[length=2mm]}, dashed, draw=grisTexte},
  actor/.pic={
    \draw[thick] (0,0.55) circle (0.17);
    \draw[thick] (0,0.38) -- (0,-0.12) (-0.3,0.22) -- (0.3,0.22) (0,-0.12) -- (-0.22,-0.5) (0,-0.12) -- (0.22,-0.5);}]
  % acteurs
  \pic at (-1.2,0.2) {actor}; \node[align=center] at (-1.2,-0.75) {Administrateur\\(opérateur)};
  \pic at (14.4,1.6) {actor}; \node[align=center] at (14.4,0.65) {Agent de collecte\\«système»};
  \pic at (14.4,-3.4) {actor}; \node[align=center] at (14.4,-4.35) {Service de\\notification «système»};
  % frontière du système
  \draw[rounded corners=6pt, draw=grisTexte, thick] (0.6,4.4) rectangle (12.6,-5.4);
  \node[font=\sffamily\footnotesize\bfseries, anchor=north west] at (0.8,4.3) {SentinelOps};
  % cas utilisateur
  \node[uc] (auth) at (3.4,3.4) {S'authentifier};
  \node[uc] (ent)  at (3.4,2.2) {Consulter l'état des entités};
  \node[uc] (met)  at (3.4,1.0) {Visualiser les métriques};
  \node[uc] (inc)  at (3.4,-0.2) {Consulter les incidents};
  \node[uc] (ack)  at (3.4,-1.4) {Acquitter / résoudre une alerte};
  \node[uc] (rul)  at (3.4,-2.6) {Gérer les règles};
  \node[uc] (set)  at (3.4,-3.8) {Configurer les notifications};
  \node[uc] (ops)  at (3.4,-5.0) {Consulter les journaux /\\redémarrer un conteneur};
  \node[uc] (det)  at (8.8,-0.2) {Consulter le détail d'une alerte};
  \node[uc] (pm)   at (8.8,-1.5) {Exporter un post-mortem};
  % cas système
  \node[uc] (col)  at (9.6,3.4) {Collecter les métriques};
  \node[uc] (ano)  at (9.6,2.2) {Détecter seuils et anomalies};
  \node[uc] (cyc)  at (9.6,1.0) {Déclencher / résoudre les alertes};
  \node[uc] (snd)  at (9.6,-3.4) {Envoyer les notifications};
  % liens acteurs
  \foreach \u in {auth,ent,met,inc,ack,rul,set,ops} \draw[lk] (-0.85,0) -- (\u.west);
  \foreach \u in {col,ano,cyc} \draw[lk] (14.0,1.4) -- (\u.east);
  \draw[lk] (14.0,-3.6) -- (snd.east);
  % relations
  \draw[inc] (inc.east) -- node[above, font=\sffamily\scriptsize, text=grisTexte]{«extend»} (det.west);
  \draw[inc] (pm.north) -- node[right, font=\sffamily\scriptsize, text=grisTexte]{«extend»} (det.south);
\end{tikzpicture}
}
\caption{Diagramme de cas d'utilisation global}
\label{fig:cu-global}
\end{figure}

\section{Description des cas d'utilisation principaux}

\newcommand{\fichecu}[6]{%
\begin{table}[htbp]\centering\small
\begin{tabularx}{\linewidth}{L{3.2cm}Y}
\toprule
\multicolumn{2}{l}{\textbf{Cas d'utilisation : #1}}\\
\midrule
Acteur & #2 \\
Préconditions & #3 \\
Scénario nominal & #4 \\
Alternatives & #5 \\
Postconditions & #6 \\
\bottomrule
\end{tabularx}
\caption{Description du cas \og #1 \fg{}}
\end{table}}

\fichecu{S'authentifier}{Administrateur}
{la plateforme est démarrée.}
{1.~L'administrateur saisit son identifiant et son mot de passe. 2.~Le système vérifie
l'empreinte du mot de passe. 3.~Le système crée une session et ouvre le tableau de bord.}
{2a.~Identifiants incorrects : message d'erreur ; après 5 échecs en 15 minutes depuis la même
adresse, la connexion est bloquée temporairement. 3a.~Première connexion : le système impose le
choix d'un nouveau mot de passe d'au moins 8 caractères avant tout accès.}
{une session valable 12 heures (2 heures d'inactivité au plus) est ouverte.}

\fichecu{Traiter un incident}{Administrateur}
{l'administrateur est connecté ; une alerte est déclenchée.}
{1.~L'incident apparaît en tête de la liste. 2.~L'administrateur ouvre le détail de l'alerte.
3.~Il consulte la cause (valeurs mesurées, courbe, processus). 4.~Il acquitte l'alerte.
5.~Une fois le problème corrigé, l'alerte se résout automatiquement et une notification de
résolution est envoyée.}
{4a.~Il résout l'alerte manuellement. 4b.~Pour un conteneur, il consulte ses journaux ou le
redémarre. 5a.~Il exporte le post-mortem au format Markdown.}
{l'alerte est résolue et son historique reste consultable 30 jours.}

\begin{conclusionchapitre}
Les besoins étant spécifiés, le chapitre suivant présente la conception retenue pour y
répondre : architecture, données, automate d'alertes et détection d'anomalies.
\end{conclusionchapitre}
